\documentclass[11pt]{article}
\usepackage[a4paper,margin=1in]{geometry}
\usepackage{amsmath,amssymb,amsthm,mathtools}
\usepackage{bm}
\usepackage{enumitem}
\usepackage{hyperref}
\usepackage{array}
\usepackage{booktabs}
\usepackage{tabularx}

\title{Theoretical Analysis of Safety Guarantees with Shielding}
\author{}
\date{}

\newtheorem{definition}{Definition}
\newtheorem{proposition}{Proposition}
\newtheorem{theorem}{Theorem}
\newtheorem{remark}{Remark}

\begin{document}

\maketitle

\section{Overview}

The current research line studies a \emph{layered shielding framework} for multi-UAV cooperative search. The key distinction is no longer whether safety is turned on or off, but rather which level of conservativeness is active:
\begin{itemize}[leftmargin=1.5em]
    \item \textbf{Hard-safety layer}: one-step safety is always enforced before execution.
    \item \textbf{Recursive-feasibility layer}: the shield prefers actions that preserve at least one future safe action.
    \item \textbf{Look-ahead layer}: the shield can be extended to short-horizon viability checks.
    \item \textbf{Scheduling layer}: training progress and online risk determine how strongly the stricter layers should intervene.
\end{itemize}

Under this interpretation, the main theoretical question is not whether progressive shielding occasionally permits unsafe actions, but how hard safety, recursive feasibility, short-horizon viability, and risk-aware conservativeness can be combined into a unified and computationally tractable decision process.

\section{Safe Set and Shield Operator}

\begin{definition}[Safe state set]
Let $s_t$ denote the system state at time $t$, including UAV positions, threat positions, and other necessary environmental variables. The safe state set is defined as
\begin{equation}
\mathcal S_{\mathrm{safe}}
=
\left\{
s_t \,\middle|\,
\|u_{i,t}-u_{j,t}\| \ge R_U^{\min},\ \forall i \neq j;\ 
\|u_{i,t}-e_{k,t}\| \ge R_E^{\min},\ \forall i,k
\right\},
\end{equation}
where $R_U^{\min}$ is the minimum safety distance between any two UAVs, and $R_E^{\min}$ is the minimum safety distance between a UAV and a threat.
\end{definition}

\begin{definition}[Safe joint action set]
Let $A$ denote the joint action space and let $f(s_t,a_t)$ denote the state transition function. The safe joint action set at state $s_t$ is defined as
\begin{equation}
A^{\mathrm{safe}}(s_t)
=
\left\{
a_t \in A \mid f(s_t,a_t) \in \mathcal S_{\mathrm{safe}}
\right\}.
\end{equation}
In other words, a joint action is safe if the next state remains inside the safe state set after applying that action.
\end{definition}

\begin{definition}[Shield operator]
Define the shield operator as an \emph{action-set filter} rather than a hard-coded external action selector. Given the current state $s_t$, the shield produces an allowed action set
\begin{equation}
\mathcal A_t^{\mathrm{allow}}(s_t) \subseteq A^{\mathrm{safe}}(s_t).
\end{equation}
The actor then re-selects an action from this set through a masked policy
\begin{equation}
\pi_\theta^{\mathrm{mask}}(a \mid o_t,\mathcal A_t^{\mathrm{allow}})
=
\frac{\pi_\theta(a\mid o_t)\mathbf 1[a\in \mathcal A_t^{\mathrm{allow}}]}
{\sum_{a' \in A}\pi_\theta(a'\mid o_t)\mathbf 1[a'\in \mathcal A_t^{\mathrm{allow}}]},
\end{equation}
and the final executed action is sampled or selected from $\pi_\theta^{\mathrm{mask}}$. This formulation preserves actor-side preference learning while ensuring that only allowed actions can be executed.
\end{definition}

\begin{definition}[One-step recursive-feasible action set]
The one-step recursive-feasible action set is defined as
\begin{equation}
A^{\mathrm{rec}}(s_t)
=
\left\{
 a_t \in A^{\mathrm{safe}}(s_t)
 \mid
 A^{\mathrm{safe}}(f(s_t,a_t)) \neq \varnothing
\right\}.
\end{equation}
An action is recursive-feasible if it is safe now and does not immediately lead to a state with no remaining safe action.
\end{definition}

\begin{definition}[Small-horizon viable action set]
For a finite horizon $H \ge 1$, define the small-horizon viable action set
\begin{equation}
A_H^{\mathrm{viable}}(s_t)
=
\left\{
 a_t \in A
 \mid
 \exists \text{ a safe continuation of length } H \text{ after } a_t
\right\}.
\end{equation}
This set generalizes one-step recursive feasibility to short-horizon look-ahead safety. In practice, the current implementation may use a greedy or approximate future-safety existence test as a tractable surrogate.
\end{definition}

\section{Exact Joint Modeling of the One-Step Hard-Safety Layer}

The current implementation computes a sequential approximation of the one-step hard-safe action set. To clarify the underlying mathematical object and to diagnose whether observed dead-end states are intrinsic or induced by the sequential approximation, it is useful to define the exact one-step hard-safety problem as a joint multi-constraint feasibility model.

\begin{definition}[Per-agent valid action domain]
For each UAV $i$, let
\begin{equation}
D_i(s_t)
=
\left\{
a_i \,\middle|\,
\text{$a_i$ satisfies the local kinematic and boundary-validity constraints at state $s_t$}
\right\}.
\end{equation}
Equivalently, if UAV $i$ is currently at $(x_{i,t},y_{i,t},z_{i,t})$ and action $a_i$ induces displacement $(\Delta x_{i,a_i},\Delta y_{i,a_i},\Delta z_{i,a_i})$, then
\begin{equation}
0 \le x_{i,t} + \Delta x_{i,a_i} \le S_x-1,
\end{equation}
\begin{equation}
0 \le y_{i,t} + \Delta y_{i,a_i} \le S_y-1,
\end{equation}
\begin{equation}
0 \le z_{i,t} + \Delta z_{i,a_i} \le Z-1,
\end{equation}
must hold, together with the platform's local turning and motion constraints. Thus, the boundary restriction can be treated either as an explicit inequality in the optimization model or as a preprocessing rule that removes invalid actions from $D_i(s_t)$.
\end{definition}

\begin{definition}[Exact joint one-step hard-safe feasibility model]
Introduce binary decision variables
\begin{equation}
x_{i,a} \in \{0,1\},
\end{equation}
where $x_{i,a}=1$ means that UAV $i$ executes action $a \in D_i(s_t)$. The exact one-step hard-safe model is the set of all binary assignments satisfying the following constraints:
\begin{align}
\sum_{a \in D_i(s_t)} x_{i,a} &= 1, && \forall i, \label{eq:joint-hard-one-action}\\
x_{i,a} &= 0, && \forall a \in T_i(s_t), \label{eq:joint-hard-threat}\\
x_{i,a} + x_{j,b} &\le 1, && \forall ((i,a),(j,b)) \in C_{ij}^{\mathrm{dist}}(s_t), \label{eq:joint-hard-pair}\\
x_{i,a} + x_{j,b} &\le 1, && \forall ((i,a),(j,b)) \in C_{ij}^{\mathrm{swap}}(s_t). \label{eq:joint-hard-swap}
\end{align}
Here $T_i(s_t)$ denotes the set of actions whose next-step position violates the UAV-threat hard-safety constraint, $C_{ij}^{\mathrm{dist}}(s_t)$ denotes the set of pairwise action combinations that violate the UAV-UAV distance constraint, and $C_{ij}^{\mathrm{swap}}(s_t)$ denotes the set of pairwise action combinations that induce a forbidden position swap. The exact joint one-step hard-safe feasible set is therefore
\begin{equation}
\mathcal X_{\mathrm{hard}}(s_t)
=
\left\{
x \,\middle|\,
\text{\eqref{eq:joint-hard-one-action}--\eqref{eq:joint-hard-swap} hold}
\right\}.
\end{equation}
\end{definition}

\begin{remark}
This model is a pure feasibility problem and does not require an optimization objective in order to define hard safety. In practice, a lightweight objective aligned with actor preference can optionally be added only to extract a witness completion, but the shield semantics should remain unchanged: the actor proposes first, the shield constructs an allowed action set, and the actor re-selects inside that set if needed.
\end{remark}

\begin{definition}[Exact projected one-step hard-safe action set]
The exact admissible one-step hard-safe set for UAV $i$ is obtained by projection of the joint feasible set:
\begin{equation}
A^{\star}_{\mathrm{hard},i}(s_t)
=
\left\{
a \in D_i(s_t)
\,\middle|\,
\exists x \in \mathcal X_{\mathrm{hard}}(s_t)\ \text{such that}\ x_{i,a}=1
\right\}.
\end{equation}
In words, action $a$ is admissible for UAV $i$ if and only if there exists a jointly hard-safe completion of all other UAV actions under which $i$ can execute $a$.
\end{definition}

\begin{definition}[Conditional exact hard-safe action set]
Suppose a subset of UAVs $\mathcal P \subseteq \{1,\dots,N\}$ has already been assigned actions $\hat a_{\mathcal P} = \{\hat a_k\}_{k \in \mathcal P}$. The corresponding conditional joint feasible set is
\begin{equation}
\mathcal X_{\mathrm{hard}}(s_t;\mathcal P,\hat a_{\mathcal P})
=
\left\{
x \in \mathcal X_{\mathrm{hard}}(s_t)
\,\middle|\,
x_{k,\hat a_k}=1,\ \forall k \in \mathcal P
\right\}.
\end{equation}
Then the conditional admissible set for UAV $i \notin \mathcal P$ is
\begin{equation}
A^{\star}_{\mathrm{hard},i}(s_t;\mathcal P,\hat a_{\mathcal P})
=
\left\{
a \in D_i(s_t)
\middle|\,
\exists x \in \mathcal X_{\mathrm{hard}}(s_t;\mathcal P \cup \{i\},\hat a_{\mathcal P} \cup \{a\})
\right\}.
\end{equation}
This conditional definition matches the actor--shield interaction semantics more closely than a naive sequential approximation based on fixing all remaining UAVs to a single reference joint action.
\end{definition}

\begin{remark}
The sequential shield currently used in implementation can be interpreted as an approximation to $A^{\star}_{\mathrm{hard},i}(s_t;\mathcal P,\hat a_{\mathcal P})$. This interpretation is useful because it separates two different causes of a dead-end:
\begin{itemize}[leftmargin=1.5em]
    \item a \textbf{true dead-end}, where the conditional exact set is genuinely empty; and
    \item an \textbf{approximation-induced dead-end}, where the exact conditional set is nonempty but the sequential approximation returns an empty set.
\end{itemize}
This distinction provides a principled diagnostic target for future algorithmic improvement.
\end{remark}

\section{Relation to MPC-like Safety Filtering}

The exact one-step hard-safety model above is closely related to a short-horizon constrained planning problem. In particular, if one only wants a \emph{single} jointly safe action at the current state, a natural MPC-like formulation is
\begin{equation}
a_t^{\mathrm{mpc}}
\in
\arg\max_{a_t \in A}
J(s_t,a_t)
\quad
\text{s.t.}
\quad
f(s_t,a_t)\in \mathcal S_{\mathrm{safe}},
\end{equation}
or, for a finite horizon $H \ge 1$,
\begin{equation}
\mathbf a_{t:t+H-1}^{\mathrm{mpc}}
\in
\arg\max_{\mathbf a_{t:t+H-1}}
J_H(s_t,\mathbf a_{t:t+H-1})
\quad
\text{s.t.}
\quad
s_{\tau+1}=f(s_\tau,a_\tau),\ 
s_{\tau+1}\in \mathcal S_{\mathrm{safe}}
\end{equation}
for all $\tau=t,\dots,t+H-1$.

Therefore, the current shielded framework is indeed \emph{MPC-like} in the following sense:
\begin{itemize}[leftmargin=1.5em]
    \item it uses an explicit dynamics or transition model;
    \item it enforces boundary, threat, collision, and swap constraints;
    \item it reasons over one-step or short-horizon future feasibility;
    \item it can use joint feasibility search to verify whether a candidate action admits a safe completion.
\end{itemize}

However, the proposed framework should \emph{not} be interpreted as a standard MPC controller. The key semantic difference is that MPC usually returns a \emph{single optimized action}, whereas the shield computes an \emph{allowed action set}. More precisely, for UAV $i$ the shield is interested in the projected feasible set
\begin{equation}
A^{\star}_{\mathrm{hard},i}(s_t)
=
\left\{
a_i \in D_i(s_t)
\middle|\,
\exists a_{-i}\ \text{s.t.}\ (a_i,a_{-i}) \in A^{\mathrm{safe}}(s_t)
\right\},
\end{equation}
or its stronger recursive / look-ahead counterparts, rather than in a single optimizer of a horizon cost.

The execution semantics are therefore:
\begin{enumerate}[label=(\arabic*),leftmargin=1.8em]
    \item the actor first proposes an action according to its learned preference;
    \item the shield computes $A_{\mathrm{hard}}$, $A^{\mathrm{rec}}$, or $A_H^{\mathrm{viable}}$ as an allowed set;
    \item if the actor's proposal is not admissible, the actor re-selects \emph{inside} the allowed set.
\end{enumerate}
Thus, even when an exact joint feasibility solver or a short-horizon look-ahead search is used internally, its role is to certify or prune admissibility, not to replace the policy with a planner.

\begin{remark}
The cleanest positioning is that the current method is a \emph{policy-preserving MPC-like shielding framework}: it borrows constrained short-horizon feasibility reasoning from MPC-style methods, but it preserves the actor--shield division of labor. This distinction is central to the present research line, because progressive scheduling and risk-aware scheduling are designed to control \emph{how strongly the allowed set is shrunk}, rather than to decide when a model-predictive planner should take over execution.
\end{remark}

\section{One-Step Safety Guarantee}
 
\begin{proposition}[One-step safety]
Suppose the current state satisfies $s_t \in \mathcal S_{\mathrm{safe}}$, and the allowed action set returned by the shield satisfies
\begin{equation}
\mathcal A_t^{\mathrm{allow}}(s_t) \subseteq A^{\mathrm{safe}}(s_t),
\end{equation}
with $\mathcal A_t^{\mathrm{allow}}(s_t) \neq \varnothing$. If the final executed action is re-selected by the actor from $\mathcal A_t^{\mathrm{allow}}(s_t)$, then the next state satisfies
\begin{equation}
s_{t+1} = f(s_t,\tilde a_t) \in \mathcal S_{\mathrm{safe}}.
\end{equation}
\end{proposition}

\begin{proof}
Because $\mathcal A_t^{\mathrm{allow}}(s_t) \subseteq A^{\mathrm{safe}}(s_t)$, any action $\tilde a_t$ selected from $\mathcal A_t^{\mathrm{allow}}(s_t)$ must also belong to $A^{\mathrm{safe}}(s_t)$. By definition of $A^{\mathrm{safe}}(s_t)$, this implies
\begin{equation}
f(s_t,\tilde a_t) \in \mathcal S_{\mathrm{safe}}.
\end{equation}
Therefore, $s_{t+1} \in \mathcal S_{\mathrm{safe}}$.
\end{proof}

\section{Recursive Safety Guarantee}

\begin{theorem}[Recursive safety]
Assume that:
\begin{enumerate}[label=(\arabic*),leftmargin=1.8em]
    \item the initial state is safe, i.e., $s_0 \in \mathcal S_{\mathrm{safe}}$;
    \item for any $s_t \in \mathcal S_{\mathrm{safe}}$, the shield returns a nonempty allowed set satisfying
    \begin{equation}
    \mathcal A_t^{\mathrm{allow}}(s_t) \subseteq A^{\mathrm{safe}}(s_t),
    \end{equation}
    \item the actor always re-selects the final action from $\mathcal A_t^{\mathrm{allow}}(s_t)$;
    \item the shield uses a transition model consistent with the actual environment transition.
\end{enumerate}
Then, for all $t \ge 0$,
\begin{equation}
s_t \in \mathcal S_{\mathrm{safe}}.
\end{equation}
\end{theorem}

\begin{proof}
We use mathematical induction.

\textbf{Base case.} By Assumption (1), $s_0 \in \mathcal S_{\mathrm{safe}}$.

\textbf{Induction step.} Assume that $s_t \in \mathcal S_{\mathrm{safe}}$ for some $t \ge 0$. By Assumptions (2) and (3), the final action $\tilde a_t$ belongs to $A^{\mathrm{safe}}(s_t)$. By Proposition 1, this implies that
\begin{equation}
s_{t+1} = f(s_t,\tilde a_t) \in \mathcal S_{\mathrm{safe}}.
\end{equation}
Hence the claim holds for $t+1$.

By induction, $s_t \in \mathcal S_{\mathrm{safe}}$ for all $t \ge 0$.
\end{proof}

\section{Robust Safety under Bounded Prediction Errors}

In practical systems, the shield may rely on predicted next-step UAV and threat positions. Therefore, safety guarantees should account for prediction errors.

\begin{theorem}[Robust safety under bounded errors]
Suppose the prediction errors satisfy
\begin{equation}
\|\hat u_{i,t+1} - u_{i,t+1}\| \le \epsilon_U,\qquad
\|\hat e_{k,t+1} - e_{k,t+1}\| \le \epsilon_E,
\end{equation}
where $\hat u_{i,t+1}$ and $\hat e_{k,t+1}$ are predicted UAV and threat positions, while $u_{i,t+1}$ and $e_{k,t+1}$ are the corresponding true positions.

Assume the shield enforces the tightened constraints
\begin{equation}
\|\hat u_{i,t+1} - \hat u_{j,t+1}\| \ge R_U^{\min} + 2\epsilon_U,
\end{equation}
\begin{equation}
\|\hat u_{i,t+1} - \hat e_{k,t+1}\| \ge R_E^{\min} + \epsilon_U + \epsilon_E.
\end{equation}
Then the true next-step positions satisfy the original safety constraints
\begin{equation}
\|u_{i,t+1} - u_{j,t+1}\| \ge R_U^{\min},
\end{equation}
\begin{equation}
\|u_{i,t+1} - e_{k,t+1}\| \ge R_E^{\min}.
\end{equation}
\end{theorem}

\begin{proof}
By the triangle inequality,
\begin{align}
\|u_{i,t+1} - u_{j,t+1}\|
&\ge
\|\hat u_{i,t+1} - \hat u_{j,t+1}\|
- \|u_{i,t+1} - \hat u_{i,t+1}\|
- \|u_{j,t+1} - \hat u_{j,t+1}\| \\
&\ge
(R_U^{\min} + 2\epsilon_U) - \epsilon_U - \epsilon_U \\
&=
R_U^{\min}.
\end{align}
Similarly,
\begin{align}
\|u_{i,t+1} - e_{k,t+1}\|
&\ge
\|\hat u_{i,t+1} - \hat e_{k,t+1}\|
- \|u_{i,t+1} - \hat u_{i,t+1}\|
- \|e_{k,t+1} - \hat e_{k,t+1}\| \\
&\ge
(R_E^{\min} + \epsilon_U + \epsilon_E) - \epsilon_U - \epsilon_E \\
&=
R_E^{\min}.
\end{align}
Therefore, the original safety constraints still hold under bounded prediction errors.
\end{proof}

\section{Progressive Conservativeness during Training}

The current research line does not interpret progressive shielding as occasionally allowing unsafe actions to pass through. Instead, one-step hard safety is always enforced, while the \emph{conservativeness level} of the shield is gradually strengthened over training.

A convenient hierarchy is
\begin{equation}
\mathcal A^{(0)}(s_t)=A^{\mathrm{safe}}(s_t),
\qquad
\mathcal A^{(1)}(s_t)=A^{\mathrm{rec}}(s_t),
\qquad
\mathcal A^{(2)}(s_t)=A_H^{\mathrm{viable}}(s_t).
\end{equation}
A time-varying conservativeness index $\kappa(t)$ selects the active allowed action set:
\begin{equation}
\mathcal A_t^{\mathrm{allow}}(s_t)=\mathcal A^{(\kappa(t))}(s_t).
\end{equation}
Thus, the progressive mechanism upgrades the decision process from one-step safety to recursive feasibility and then to short-horizon viability, without relaxing the hard safety constraint itself.

When the actor's proposed action does not belong to the current allowed action set, the shield is triggered and the actor re-selects within $\mathcal A_t^{\mathrm{allow}}(s_t)$. A trigger-based penalty can then be added to the reward:
\begin{equation}
r_t' = r_t - \beta_t \, \mathbb I_{\mathrm{shield}}(t).
\end{equation}
This makes progressive shielding compatible with policy learning: the shield restricts unsafe or overly aggressive decisions, while the actor is still responsible for choosing within the permitted action set.

\begin{remark}
Under this interpretation, progressive shielding is a \emph{progressive conservativeness} mechanism rather than a relaxation of hard safety. Its role is to control how strongly recursive-feasible or look-ahead constraints influence the actor over the course of training.
\end{remark}


\section{Risk-Aware Dual Scheduling}

The next step is not to directly train a learned risk network, but to upgrade the current binary high-risk gate into a \emph{continuous, interpretable, and low-cost} risk score. The first formal version is
\begin{equation}
\xi_{i,t} = w_1 \, \xi_{i,t}^{\mathrm{clear}} + w_2 \, \xi_{i,t}^{\mathrm{region}} + w_3 \, \xi_{i,t}^{\mathrm{hist}},
\qquad
w_1+w_2+w_3=1,
\quad w_k\ge 0.
\end{equation}
A practical initialization is $w_1=0.5$, $w_2=0.3$, and $w_3=0.2$, so that geometric clearance is treated as the dominant term.

The three components are defined as follows.

\paragraph{Clearance risk.}
Let $m_{i,t}$ denote the current minimum candidate safety clearance (the same quantity used by the existing rule-based gate). Define
\begin{equation}
\xi_{i,t}^{\mathrm{clear}}
=
\mathrm{clip}\left(1-\frac{m_{i,t}}{M_c},0,1\right),
\end{equation}
where $M_c$ is a normalization constant.

\paragraph{Region risk.}
This term captures whether the agent is currently located in a locally risky area, such as near the boundary, close to threats, or inside a crowded UAV neighborhood. A simple form is
\begin{equation}
\xi_{i,t}^{\mathrm{region}}
=
\frac{\mathbb I_{\mathrm{boundary}} + \tilde n^{\mathrm{threat}}_{i,t} + \mathbb I_{\mathrm{crowded}}}{3},
\end{equation}
where $\tilde n^{\mathrm{threat}}_{i,t}$ is the normalized local threat count.

\paragraph{History risk.}
Let $W$ be a recent time window. Define the trigger-frequency term
\begin{equation}
\xi_{i,t}^{\mathrm{hist}}
=
\frac{1}{W} \sum_{\tau=t-W+1}^{t} \mathbb I_{\mathrm{shield}}(i,\tau),
\end{equation}
which reflects whether the current agent has been repeatedly corrected by the shield in the recent past.

The key motivation for using these three terms first is computational: they are already available or cheaply computable from the current shield pipeline, so they do not require an additional expensive online enumeration of $A^{\mathrm{safe}}$.

At the current stage, the risk score is primarily used to decide whether an agent should remain in one-step hard-safe mode or be upgraded to recursive-feasible checking. A simple scheduling rule is
\begin{equation}
\xi_{i,t} > \eta
\quad \Rightarrow \quad
\text{activate stronger shield logic for agent } i.
\end{equation}
This keeps the first risk-aware design tightly coupled to the current implementation.

\paragraph{TODO items for later extensions.}
Two additional terms are intentionally postponed:
\begin{itemize}[leftmargin=1.5em]
    \item \textbf{Feasibility-proxy risk:}
    \begin{equation}
    \xi_{i,t}^{\mathrm{feas\mbox{-}proxy}} = 1 - \frac{|A_i^{\mathrm{rule}}(s_t)|}{|A_i|},
    \end{equation}
    or a top-$k$ survival-ratio variant.
    \item \textbf{Uncertainty / preference-conflict risk:} based on whether the actor's preferred actions are rejected by the rule mask, or on a small action-gap / logit-gap indicator.
\end{itemize}
These terms are left as explicit TODOs because the first goal is to validate that the low-cost three-term score already provides useful scheduling behavior.

\section{Recommended Theoretical Positioning}

For the paper, the cleanest theoretical positioning is:
\begin{itemize}[leftmargin=1.5em]
    \item \textbf{Hard safety layer}: one-step safety is always enforced through action-set filtering and actor-constrained re-selection.
    \item \textbf{Conservativeness layer}: recursive-feasible and small-horizon look-ahead sets are used to strengthen future viability when needed.
    \item \textbf{Scheduling layer}: training progress and online risk determine how strongly the stricter shield logic should intervene.
\end{itemize}

This positioning preserves a clear separation of roles: hard safety provides the basic guarantee, recursive-feasible and look-ahead shielding improve sustainable safety, and progressive/risk-aware scheduling controls how much conservativeness is injected during learning.

\section{Summary}

The main conclusions are as follows:
\begin{enumerate}[label=(\arabic*),leftmargin=1.8em]
    \item If the final executed action is always re-selected from an allowed set contained in $A^{\mathrm{safe}}(s_t)$, one-step safety can be guaranteed.
    \item If the initial state is safe and the shield always returns a nonempty allowed set contained in $A^{\mathrm{safe}}(s_t)$, recursive safety follows by induction.
    \item If prediction errors are bounded, robust safety can be recovered by tightening the shield constraints.
    \item One-step recursive-feasible shielding and small-horizon look-ahead shielding strengthen \emph{sustainable safety} by avoiding immediate transitions into dead-end states.
    \item Progressive and risk-aware scheduling should be interpreted as mechanisms for controlling conservativeness, intervention frequency, and action-space shrinkage --- not as permissions to violate hard safety.
\end{enumerate}

Therefore, the current research line is best summarized as a layered framework: one-step hard safety is always enforced, recursive-feasible and short-horizon viable sets improve future feasibility, and progressive/risk-aware scheduling determines how aggressively these stronger filters shape learning and execution.

\section{Current Research Roadmap}

The current plan is to keep the original three-point research line rather than switching the main contribution to learned safe masks. Concretely, the next scientific stages are:
\begin{enumerate}[label=(\arabic*),leftmargin=1.8em]
    \item establish formal baselines for \texttt{off}, \texttt{safe}, and \texttt{recursive};
    \item formalize the current rule-based gate as a continuous risk function using the low-cost \emph{clear + region + hist} structure;
    \item use this risk function to decide when the shield should escalate from one-step hard safety to recursive-feasible checking, and then replace the current greedy future-safe existence approximation with an $H=2$ small-horizon look-ahead shield;
    \item introduce progressive conservativeness so that the shield evolves from one-step safety to recursive-feasible and then short-horizon viable filtering over training;
    \item finally extend the system to full dual scheduling of intervention strength and action-space shrinkage.
\end{enumerate}

An important implementation lesson is that exact recursive shielding should not be naively expanded into full joint future-action enumeration at every step. However, this computational concern refines the \emph{implementation strategy}, not the main scientific direction. The main contribution remains the unified combination of progressive shielding, recursive-feasible / short-horizon viability maintenance, and risk-aware dual scheduling.

\section{Comparison with Related Shielding Directions}

Table~\ref{tab:related-work-comparison} summarizes the positioning of the proposed research direction against the most relevant shielding-related lines of work discussed during the literature scan. The goal is not to claim that shielding, look-ahead, or adaptivity are entirely new, but to clarify that the novelty should be placed on the \emph{unified combination} of training-progress-aware shielding, recursive-feasibility preservation, and risk-aware dual scheduling in the multi-UAV cooperative search setting.

\begin{table}[htbp]
\centering
\caption{Comparison between the planned framework and closely related shielding directions}
\label{tab:related-work-comparison}
\renewcommand{\arraystretch}{1.15}
\begin{tabularx}{\textwidth}{>{\raggedright\arraybackslash}p{3.0cm} >{\raggedright\arraybackslash}p{3.1cm} >{\raggedright\arraybackslash}p{2.2cm} >{\raggedright\arraybackslash}p{2.2cm} >{\raggedright\arraybackslash}X}
\toprule
\textbf{Direction / Representative Work} & \textbf{Main Idea} & \textbf{Look-ahead Shield} & \textbf{Adaptive / Progressive Control} & \textbf{Main Difference from Planned Framework} \\
\midrule
Classical shielded RL & Use a shield to block unsafe actions before execution & Usually no & Usually fixed & Establishes the basic shielding idea, but does not address training-progress-aware intervention, recursive feasibility preservation, or multi-UAV search coupling. \\

Dynamic multi-agent shielding \newline (e.g., model-based dynamic shielding for safe and efficient MARL, AAMAS 2023) & Dynamically construct/update shields in multi-agent settings; includes model-based and $k$-step look-ahead elements & Yes & Dynamic shield construction & This is the closest line. The planned framework should distinguish itself by emphasizing training-progress-aware intervention strength, action-space shrinkage scheduling, and recursive-feasibility-oriented control for multi-UAV target search. \\

Decentralized / shield decentralization \newline (e.g., NeurIPS 2022) & Decompose or decentralize shielding for multi-agent execution & Not the main focus & Limited & Focuses on decentralized realizability of shields rather than progressive training control or feasibility-aware action-space shrinkage. \\

POMDP / shield decomposition \newline (e.g., RLC 2024) & Extend shielding ideas to partially observable multi-agent environments & Not the main focus & Limited & Shares the partial-observability setting, but does not center on progressive scheduling or the preservation of future safe action availability. \\

Adaptive shielding \newline (e.g., ADVICE, 2025) & Adjust shielding behavior based on learned safe/unsafe knowledge or black-box interaction & Partial & Yes & Related in spirit to adaptive safety intervention, but not specifically organized around dual scheduling of intervention strength and action-space conservativeness in cooperative UAV search. \\

Adaptive robust MPC shielding \newline (e.g., adaptive robust model predictive shielding, 2025) & Use robust/model-predictive ideas to adapt safety margins or shielding conservativeness & Yes & Yes & Relevant to adaptive conservativeness, but more MPC/process-control oriented than MARL-oriented, and not directly focused on recursive-feasibility-preserving action pruning in multi-UAV search. \\

Planned framework \newline (this work) & Unify progressive shielding, one-step recursive-feasible shielding, small-horizon look-ahead shielding, and risk-aware dual scheduling & Yes & Yes & Targets the exploration-safety-feasibility tradeoff jointly: early-stage exploration, late-stage safety strengthening, dead-end avoidance, and adaptive action-space shrinkage under multi-UAV cooperative search constraints. \\
\bottomrule
\end{tabularx}
\end{table}

\paragraph{Recommended positioning.}
The planned work should avoid claiming that shielding, look-ahead shielding, or adaptive shielding are individually novel. Instead, the strongest positioning is that the contribution lies in a \emph{unified framework} that jointly addresses:
\begin{itemize}[leftmargin=1.5em]
    \item training-progress-aware exploration-safety balancing,
    \item recursive-feasibility preservation to avoid dead-end states, and
    \item risk-aware dual scheduling of intervention strength and action-space shrinkage
\end{itemize}
for multi-UAV cooperative search with moving targets and threats.

\end{document}
